(1) 令高斯相位增量 \(X=\phi(\boldsymbol r+\boldsymbol\rho)-\phi(\boldsymbol r)\)，有 \(\langle X\rangle=0\)、\(\operatorname{Var}(X)=D_\phi(\rho)\)。零均值高斯的特征函数 \(\langle e^{iX}\rangle=\exp[-\operatorname{Var}(X)/2]\)，故
\[
\boxed{\Gamma(\rho)=\exp[-D_\phi(\rho)/2]
=\exp[-3.44(\rho/r_0)^{5/3}]}.
\]
该结果使用了高斯增量而非任意随机相位的普适恒等式；若结构函数相同但高阶累积量不为零，互相干函数可不同。

(2) 长曝光光学传递函数由瞳面相差两点的互相关决定，瞳面横向间距为 \(\rho=\lambda|\boldsymbol f|\)，其中 \(\boldsymbol f\) 是焦平面角坐标共轭的空间频率（在相应成像约定下）。大气因子因此
\[
\boxed{\operatorname{OTF}_{\rm atm}(\boldsymbol f)
=\exp[-3.44(\lambda|\boldsymbol f|/r_0)^{5/3}]}.
\]
若相位屏统计均匀且与光学瞳孔不相关，系统长曝光OTF为 \(\operatorname{OTF}_{\rm aperture}\operatorname{OTF}_{\rm atm}\)。不可把理想衍射瞳孔因子删除后便称大气决定所有空间频率，即使望远镜口径小于 \(r_0\) 也需保留有限口径衍射。

(3) 给定相同湍流路径与空气统计特性，结构函数系数 \(D_\phi\propto\lambda^{-2}\rho^{5/3}\)。而参数定义给 \(D_\phi\propto r_0^{-5/3}\rho^{5/3}\)，故 \(r_0^{-5/3}\propto\lambda^{-2}\)，即
\[
\boxed{r_0\propto\lambda^{6/5}}.
\]
大口径湍流限制的角宽尺度约 \(\theta_{\rm seeing}\sim\lambda/r_0\propto\lambda^{-1/5}\)，与理想小口径衍射宽 \(\lambda/D\) 对波长的线性标度不同。具体FWHM系数依赖长曝光PSF、惯性区和外尺度假设。

(4) \(\rho\to0\) 时 \(D_\phi\to0,\Gamma\to1\)；\(\rho/r_0\to\infty\) 时 \(\Gamma\to0\)，表明长间距瞳点相位完全失相关。若湍流具有有限外尺度，结构函数在大 \(\rho\) 时可能饱和，不会严格按 \(5/3\) 幂无限增长；振幅闪烁使互相干还包含幅度协方差；非高斯相位分布需使用完整特征函数。因此本题不能替代真实大气自适应光学误差预算。
